\documentclass[11pt]{article}
\usepackage{amsmath,amssymb,amsthm}
\usepackage[margin=1in]{geometry}
\newtheorem{theorem}{Theorem}
\newtheorem{lemma}[theorem]{Lemma}
\newtheorem{proposition}[theorem]{Proposition}
\newtheorem{corollary}[theorem]{Corollary}
\theoremstyle{remark}
\newtheorem{remark}[theorem]{Remark}
\newcommand{\Q}{\mathbb Q}
\newcommand{\Z}{\mathbb Z}
\newcommand{\C}{\mathbb C}
\newcommand{\e}{\mathrm e}
\title{P4 v2: a global lower bound on $n(r,U)$ from summing thm:arc's local pole rate}
\date{2026-09-26}
\begin{document}
\maketitle

\noindent Labels: PROVED means a complete argument, given the inputs it cites (here, given
\texttt{thm:arc}/\texttt{cor:arcUV} of paper2a, which are themselves conditional on
Theorems~\texttt{thm:model} and~\texttt{thm:RJ}). VERIFIED means high-precision floating point
checked against paper2a's own reported numbers, not a fresh interval certificate. HEURISTIC means an
argument without full error control. This note continues \texttt{P4/P4\_note.tex} (``P4v1''), which
identified two missing ingredients for a Nevanlinna-theoretic resolution of P4 (is $U$ D-algebraic):
(a) a quantitative global pole-counting rate $n(r,U)$ as $r\to1^-$, summed from the local rate
\texttt{thm:arc} gives near each fixed root of unity; (b) a disc-analogue of the
Bank--Kaufman/Steinmetz counting theorem. P4v1 explicitly did not carry out the summation in (a). This
note carries it out.

\section*{Correction (post-review, Reviewer AN): the headline claim is NOT established}

\textbf{The local-rate formula \eqref{eq:localrate} below, and everything built on it
(\S\ref{sec:invert}--\ref{sec:diagonal}, Proposition~\ref{prop:main}), is wrong as a claim about
general $\zeta$.} The formula $t^\ast_j=(L+i\pi)/(2j+3/2)+O(t)$ was copied from
\texttt{prop:minusone} (paper2a, the \emph{special case} $\zeta=-1$), not from the general theorem
that actually backs \texttt{thm:arc} (\texttt{P1f\_lemma.tex}, the general saddle/Rouch\'e theorem
with $\Delta V=(\pi is/N)(L_n+i\pi s/N)$, $s\in\{1,2\}$). The general $\Delta V$ carries an extra
factor of order $1/N$ that the $\zeta=-1$ formula (a fixed small case, $N=2$) hides, so the correct
leading behavior is $1-|q^\ast_j(\zeta)|\sim C(\zeta,N)/j$ with $C(\zeta,N)=O(1/N)$, \textbf{not}
$\ell(\zeta)/(2j)$. Checking $\zeta=i$ (\texttt{prop:qi}) independently gives yet a third, different
offset ($2j+1$, from $m=-j-1$), confirming the coefficient is genuinely case-specific and not the
single formula claimed below. The numerical "corroboration" at $\zeta=i$ is correspondingly not
meaningful: at $j=1,2$ the wrong-for-this-case formula (with offset $3/2$) happens to fit slightly
\emph{better} than $\zeta=i$'s own correct offset ($2j+1$, predicting ratio $5/3=1.667$, a worse
$3.8\%$ match) -- this is $O(1/j)$-correction noise, not evidence for the general shape.

\textbf{What survives, honestly:} the \emph{qualitative} diagnosis that P4v1's back-of-envelope
$\log(1/(1-r))$-per-$\zeta$ guess conflates a harmonic-sum question with a threshold-count question,
and that the correct per-$\zeta$ count is linear (not logarithmic) in $1/(1-r)$ whatever its exact
coefficient turns out to be -- and the observation that a diagonal argument over finitely many fixed
$\zeta$ at a time needs no uniform-in-$\zeta$ control, only individually-finite thresholds. Both are
modest, correct points. The specific formula \eqref{eq:localrate}, Lemma~\ref{lem:many}'s use of it,
and Proposition~\ref{prop:main}'s conclusion are \textbf{not established} by this note as written; a
correct version would need the general $\Delta V$-based rate re-derived from \texttt{P1f\_lemma.tex}
directly, plus (per Reviewer AN) an honest count of how many admissible $(a,N)$ contribute a
non-negligible coefficient once the true $O(1/N)$ dependence is accounted for -- not attempted here.
\textbf{P4 remains exactly as open as P4v1 left it; ingredient (a) is still not carried out.}

\section*{Summary of outcome (as originally written; superseded by the correction above)}

\textbf{Claimed result, NOT established (see correction above):} for every $A>0$,
\[
   n(r,U)\ \ge\ \frac{A}{1-r}\qquad\text{for all $r$ sufficiently close to $1$},
\]
i.e.\ $\displaystyle\liminf_{r\to1^-}(1-r)\,n(r,U)=+\infty$. The argument below is retained for
the record and because its qualitative log-vs-linear point and its diagonal-argument shape are
genuine, but its quantitative content rests on the incorrect formula \eqref{eq:localrate} and must
not be cited as a proved strengthening of P4v1.

\section{The local rate, extracted precisely}\label{sec:local}

From \texttt{paper2a.tex}, \texttt{thm:arc} (\S\ref{sec:arc} there) and its corollary
\texttt{cor:arcUV}: fix $\zeta=e^{2\pi ia/N}$, $\gcd(a,N)=1$, admissible in the sense of
\texttt{cor:arcUV} ($N\ge16$, $\frac aN\in[\frac16,\frac56]$; wider ranges are available via
P1h/P1i but are not needed below). Write
\[
   \ell=\ell(\zeta):=\log\bigl|2(1-\zeta)\bigr|,\qquad L=L(\zeta):=\operatorname{Log}\bigl(-2(1-\zeta)\bigr),
\]
so $\operatorname{Re}L=\ell$. The proof of \texttt{thm:arc} produces, for each such $\zeta$, an
explicit sequence of poles $q^\ast_j(\zeta)=\zeta\,\e^{-t^\ast_j}$ of $G^T_0$ (hence, by
\texttt{cor:arcUV}, of $U$ and $V$ after transfer through $x=\pm\sqrt{q^\ast_j}$), valid for all $j$
past some finite, $\zeta$-dependent threshold $j_0(\zeta)$ (``the index beyond which $q^\ast_j$ is a
pole $\ldots$ is not made explicit'', \texttt{rem:arcUV}), with
\[
   t^\ast_j=\frac{L+i\pi}{2j+3/2}+O(t)\qquad(j\to\infty),\qquad\text{so}\qquad
   \operatorname{Re}(t^\ast_j)=\frac{\ell}{2j+3/2}\bigl(1+o(1)\bigr).
\]
Since $|q^\ast_j|=\e^{-\operatorname{Re}(t^\ast_j)}$, and $\ell>0$ on the arc (required for the
resummation to apply at all, \texttt{thm:arc}'s statement), $\operatorname{Re}(t^\ast_j)\downarrow0$
monotonically for large $j$, i.e.\ \emph{$|q^\ast_j(\zeta)|$ increases to $1$ as $j\to\infty$}, and
\begin{equation}\label{eq:localrate}
   1-|q^\ast_j(\zeta)|=\operatorname{Re}(t^\ast_j)\bigl(1+o(1)\bigr)=\frac{\ell(\zeta)}{2j}\bigl(1+o(1)\bigr)\qquad(j\to\infty).
\end{equation}

\noindent\textbf{Numerical corroboration (VERIFIED against paper2a's own reported values, no new
computation).} At $\zeta=i$ ($N=4$, outside the $N\ge16$ arc but governed by the same mechanism,
Corollary~\texttt{cor:Ui}), paper2a reports $q^\ast_1=0.1404+0.8293i$, $q^\ast_2=0.0802+0.8974i$, so
$1-|q^\ast_1|=1-0.8411=0.1589$ and $1-|q^\ast_2|=1-0.9010=0.0990$. Equation~\eqref{eq:localrate}
predicts the ratio $(1-|q^\ast_1|)/(1-|q^\ast_2|)\approx(2\cdot2+3/2)/(2\cdot1+3/2)=5.5/3.5=1.571$; the
reported values give $0.1589/0.0990=1.605$, a $2\%$ match consistent with the stated $O(t)$ correction
at these small $j$. This is a sanity check on \eqref{eq:localrate}'s shape, not a new certificate.

\section{Inverting the local rate: linear, not logarithmic, per-$\zeta$ count}\label{sec:invert}

Fix an admissible $\zeta$ and $r\in(0,1)$. Using $1-r\le-\log r\le(1-r)/r$, the two notions of
``close to the boundary'' agree to leading order as $r\to1^-$, so I write $-\log r=1-r+O((1-r)^2)$ and
suppress the difference below. By \eqref{eq:localrate}, $|q^\ast_j(\zeta)|\le r$ iff
$\operatorname{Re}(t^\ast_j)\ge-\log r$, i.e.\ (for $j$ past the threshold where \eqref{eq:localrate}
is valid)
\[
   \frac{\ell(\zeta)}{2j+3/2}(1+o(1))\ \ge\ 1-r+O((1-r)^2)
   \iff j\ \le\ \frac{\ell(\zeta)}{2(1-r)}\bigl(1+o(1)\bigr)-\frac34.
\]
So the number of $j\ge j_0(\zeta)$ with $|q^\ast_j(\zeta)|\le r$ is
\begin{equation}\label{eq:percount}
   n_\zeta(r):=\#\{j\ge j_0(\zeta):|q^\ast_j(\zeta)|\le r\}
   =\frac{\ell(\zeta)}{2}\cdot\frac1{1-r}\ -\ j_0(\zeta)\ +\ o\Bigl(\frac1{1-r}\Bigr)\qquad(r\to1^-).
\end{equation}
This is \textbf{linear} in $1/(1-r)$, with an explicit leading coefficient $\ell(\zeta)/2$, not the
$\log(1/(1-r))$ shape guessed in P4v1 \S5(a). (The $\log$ intuition would be the right one if the
question were ``how large is $\sum_j$ [something $1/j$]'' -- a harmonic-sum question -- but the
question here is ``how many terms of a sequence decaying like $1/j$ exceed a threshold $1-r$'', which
is answered by inverting $\ell/(2j)=1-r$ directly, giving a term count linear in $1/(1-r)$, not a
harmonic sum.) This is the concrete correction to P4v1's back-of-envelope estimate that carrying out
the summation actually produces.

\section{Diagonal summation over finitely many roots of unity}\label{sec:diagonal}

\begin{lemma}\label{lem:many}
There is a constant $\ell_{\min}>0$ and infinitely many pairwise distinct admissible $\zeta$ (in the
sense of \S\ref{sec:local}) with $\ell(\zeta)\ge\ell_{\min}$.
\end{lemma}
\begin{proof}
$\ell(\zeta)=\log|2(1-\zeta)|$ is continuous and strictly positive on the closed arc
$\frac aN\in[\frac16,\frac56]$ minus $\{a/N=\frac12\}$... in fact it is strictly positive and
continuous on all of $\arg\zeta\in[\frac\pi3,\frac{5\pi}3]\setminus\{0\}$, in particular on the
compact sub-arc $\arg\zeta\in[\frac\pi3+\delta,\pi-\delta]$ for any small fixed $\delta>0$, hence
bounded below there by some $\ell_{\min}(\delta)>0$. For any prime $N\ge16$ the point $a=\lfloor
N/2\rfloor$ satisfies $\gcd(a,N)=1$ automatically and $a/N\to\frac12$ as $N\to\infty$, so for $N$
large enough $a/N$ lies in $[\frac16+\delta',\frac56-\delta']$ for any fixed $\delta'>0$; there are
infinitely many primes, giving infinitely many distinct admissible $\zeta_N=e^{2\pi ia/N}$, pairwise
distinct since they have pairwise distinct arguments in $(0,2\pi)$, with $\ell(\zeta_N)$ bounded below
uniformly (by continuity, $\ell(\zeta_N)\to\ell(-1)=\log4$ as $N\to\infty$ along this family, and is
bounded below on the whole compact range visited).
\end{proof}

\begin{proposition}[main result of this note]\label{prop:main}
Assume \texttt{thm:model} and \texttt{thm:RJ} (so that \texttt{cor:arcUV} holds). Then
$\displaystyle\liminf_{r\to1^-}(1-r)\,n(r,U)=+\infty$, where $n(r,U)$ counts poles of $U$ in
$|x|\le r$ with multiplicity.
\end{proposition}
\begin{proof}
Fix $A>0$. By Lemma~\ref{lem:many}, choose finitely many distinct admissible
$\zeta_1,\ldots,\zeta_K$ (all with $N\ge16$) with $\sum_{k=1}^K\ell(\zeta_k)/2>A$; this is possible
for $K$ large enough since $\ell(\zeta_k)\ge\ell_{\min}>0$ for infinitely many of them. Each $\zeta_k$
has a finite threshold $j_0(\zeta_k)$ (\texttt{thm:arc}'s proof; not numerically given, but finite).
By \eqref{eq:percount}, for each fixed $k$,
\[
   n_{\zeta_k}(r)=\frac{\ell(\zeta_k)}2\cdot\frac1{1-r}-j_0(\zeta_k)+o\Bigl(\frac1{1-r}\Bigr)\qquad(r\to1^-).
\]
Since $K$ is fixed, summing these $K$ (finitely many) asymptotics is legitimate term-by-term:
\[
   \sum_{k=1}^Kn_{\zeta_k}(r)=\Bigl(\sum_{k=1}^K\frac{\ell(\zeta_k)}2\Bigr)\frac1{1-r}-\sum_{k=1}^Kj_0(\zeta_k)+o\Bigl(\frac1{1-r}\Bigr)
   >\;\frac{A}{1-r}\quad\text{for $r$ close enough to $1$ (depending on $K$, hence on $A$).}
\]
The poles counted at distinct $\zeta_k$ are, for $r$ close enough to $1$, genuinely distinct points
of $\{|x|\le r\}$: along the tie ray $\arg t=\theta^\ast(\zeta_k)$, $q^\ast_j(\zeta_k)\to\zeta_k$ in
\emph{both} modulus and argument as $j\to\infty$ (the perturbation from $\zeta_k$ has size
$|t^\ast_j|=O(1/j)$ in every direction, not just radially), so for fixed distinct $\zeta_1,\ldots,\zeta_K$
the corresponding pole clusters live in disjoint shrinking neighbourhoods of $\zeta_1,\ldots,\zeta_K$
once $j$ (equivalently $r$) is large enough; no double counting occurs. Transferring through
$x=\pm\sqrt{q^\ast_j}$ and complex conjugation (\texttt{cor:arcUV}) multiplies every count by an
absolute constant ($\le4$) and does not affect distinctness for the same reason (the four points
$\pm\sqrt{q^\ast_j},\pm\sqrt{\overline{q^\ast_j}}$ cluster near the four fixed points
$\pm\sqrt{\zeta_k},\pm\sqrt{\overline{\zeta_k}}$, distinct across distinct $k$ once close enough).
Hence $n(r,U)\ge\sum_{k=1}^Kn_{\zeta_k}(r)>A/(1-r)$ for $r$ close enough to $1$. Since $A>0$ was
arbitrary, $\liminf_{r\to1^-}(1-r)n(r,U)=+\infty$.
\end{proof}

\begin{remark}[why no uniform constants are needed here, and where they would be needed]\label{rem:diagonal}
The proof above needs, for each \emph{fixed} $K$, only the $K$ individual (finite, if unspecified)
thresholds $j_0(\zeta_1),\ldots,j_0(\zeta_K)$ -- a quantifier order (``for all $A$ there is $r_0(A)$'')
that lets $K$, and hence the finite list of thresholds used, depend on $A$ freely. This is why P4v1's
concern about uncontrolled implied constants does not block Proposition~\ref{prop:main}: no
uniformity \emph{across} $\zeta$ is used, only finitely-many, one-at-a-time applications of
\texttt{thm:arc}. Uniformity across $\zeta$ would only be needed to let $K=K(r)\to\infty$ \emph{as
$r\to1$} (i.e.\ to use more and more roots of unity as $r$ approaches $1$, rather than a fixed finite
list), which is exactly what would be needed to push the bound from ``faster than every fixed multiple
of $1/(1-r)$'' to ``faster than every fixed power of $1/(1-r)$'' (superpolynomial) -- see
\S\ref{sec:diagnosis}.
\end{remark}

\section{Comparison with what a disc Bank--Kaufman/Steinmetz bound would need}\label{sec:compare}

P4v1 \S3 already flags that the plane-case Malmquist/Steinmetz/Bank--Kaufman theorems bound the
Nevanlinna characteristic $T(r,y)$ (equivalently, via the first main theorem, $N(r,y)=\int_0^r
n(t,y)/t\,dt$, hence $n(r,y)$ up to the usual logarithmic-derivative slack) of a transcendental
meromorphic solution of a fixed algebraic ODE with polynomial/rational coefficients, in terms of the
order $k$ and degree $d$ of the equation and the growth of its coefficients; for coefficients of
bounded degree the resulting bound on $T(r,y)$ is of \emph{finite order} in $r$ (a fixed power of
$r$) as $r\to\infty$ on $\C$. The natural disc analogue (Tsuji-type; Bank's disc papers, per
\texttt{rem:disc} of P4v1, cited from memory and not independently re-verified here either) should
bound $T(r,y)$, as $r\to1^-$, by a fixed power of $1/(1-r)$ -- i.e.\ ``finite order'' becomes
\emph{polynomial growth in $1/(1-r)$}, for coefficients bounded by a fixed power of $1/(1-r)$ as well
(rational functions of $x$ alone have coefficients bounded near $|x|=1$ unless they have poles
exactly on the unit circle, which would need separate handling; assume for this comparison
coefficients holomorphic on $\overline D$, the case relevant to ruling out a fixed ODE with, e.g.,
polynomial coefficients).

\textbf{What Proposition~\ref{prop:main} gives, and what it does not give, against this.}
Proposition~\ref{prop:main} shows $n(r,U)$ is not $O(1/(1-r))$ for any constant -- it beats every
\emph{linear-order-1} bound in $1/(1-r)$. If the disc Bank--Kaufman analogue for a fixed order-$k$,
degree-$d$ ODE gave a bound of the shape $n(r,y)=O((1-r)^{-\rho(k,d)})$ with $\rho(k,d)=1$ for the
simplest (e.g.\ Riccati, $k=1,d=2$) case, Proposition~\ref{prop:main} would already rule out
\emph{that specific} equation class. But it does \emph{not} rule out equations with $\rho(k,d)>1$
(which presumably exist for higher $k,d$, exactly mirroring how in the plane case higher-order
algebraic ODEs admit solutions of higher finite order): to rule out \emph{every} $(k,d)$
simultaneously one needs $n(r,U)$ to beat $(1-r)^{-\rho}$ for \emph{every} fixed $\rho$, i.e.\
superpolynomial growth in $1/(1-r)$, which is a strictly stronger statement than
Proposition~\ref{prop:main} and is \textbf{not established here}.

\section{Diagnosis: what closing the superpolynomial gap needs}\label{sec:diagnosis}

Pushing \S\ref{sec:diagonal}'s argument from ``beats every constant multiple of $1/(1-r)$'' to
``beats every power of $1/(1-r)$'' means letting the number of roots of unity used, $K$, grow as
$r\to1^-$ (e.g.\ $K(r)\sim(1-r)^{-\rho}$ for growing $\rho$), and this needs, uniformly over the
$K(r)$ many $\zeta$'s used simultaneously:
\begin{enumerate}
\item A bound on $j_0(\zeta)$ (the threshold past which \eqref{eq:localrate} is valid) as a function
of $N$ (or of $\ell(\zeta)$), rather than the qualitative ``eventually true'' statement
\texttt{thm:arc}'s proof gives. \texttt{rem:arcUV} states explicitly that this index ``is not made
explicit''.
\item A bound on how many admissible $\zeta=e(a/N)$ with $N\le N_{\max}(r)$ exist, for $N_{\max}(r)$
tied to $1-r$ -- this part is elementary number theory (Euler-phi counting of reduced fractions in an
interval, $\Theta(N_{\max}^2)$ many with $N\le N_{\max}$) and is \emph{not} an obstruction.
\item Crucially, the hazard-induction certificates behind Theorems~M, M$'$, main$'$
(\texttt{P1e\_lemma.tex}, \texttt{P1g\_lemma.tex}, \texttt{P1h\_lemma.tex}) are stated as certified
for $N$ up to $10^5$ by direct computation and by ``a hand argument for $N>10^5$'' (\texttt{thm:arc}'s
proof sketch, quoting \texttt{P1h\_lemma.tex}); whether that hand argument gives an \emph{explicit,
uniform} lower bound on $|\omega|$-margins / non-vanishing constants as $N\to\infty$ (as opposed to
merely an existence statement for each large $N$) is not something this note verifies -- it would
need rereading \texttt{P1g\_lemma.tex}/\texttt{P1h\_lemma.tex} in full with this specific question
(``is the $N\to\infty$ tail bound uniform and explicit?'') in mind, which was not done here.
\end{enumerate}
If (1) and (3) can be answered with an explicit polynomial-in-$N$ (or even just \emph{any} explicit)
bound on $j_0(\zeta)$, the diagonal argument of \S\ref{sec:diagonal} upgrades mechanically: take
$K(r)\to\infty$ at a controlled rate, and the leading coefficient $\sum\ell(\zeta_k)/2$ over
$K(r)\sim(1-r)^{-\rho}$ many terms (each $\ell(\zeta_k)$ bounded below, hence the sum growing at least
linearly in $K(r)$) would give $n(r,U)\gtrsim(1-r)^{-1-\rho}\cdot(1-r)$-type improvement -- growing
without bound in $\rho$, i.e.\ superpolynomial. This is plausible but is exactly the missing
computation; it is not carried out here for lack of the explicit thresholds in (1)/(3), and is left as
the next concrete step.

\section*{Weakest points}
\begin{itemize}
\item \textbf{The whole note is conditional} on \texttt{thm:model}, \texttt{thm:RJ} (as
\texttt{cor:arcUV} already is), and on the correctness of \texttt{thm:arc}'s proof sketch as quoted
in paper2a (this note did not re-derive \texttt{thm:arc} from \texttt{P1f}/\texttt{P1h\_lemma.tex}
from scratch; it takes the stated asymptotic $t^\ast_j\sim(L+i\pi)/(2j+3/2)$ at face value from
paper2a's text).
\item \S\ref{sec:invert}'s inversion treats $1-r\approx-\log r$ and drops the stated $O(t)$/$o(1)$
correction terms uniformly; this is standard but was not made into an explicit error bound (e.g.\ no
claim of the form ``for $r>1-\epsilon_0(\zeta)$'', only ``eventually''). Making
Proposition~\ref{prop:main}'s proof fully quantitative (an explicit $r_0(A)$) would require exactly
the same threshold data flagged missing in \S\ref{sec:diagnosis}.
\item Proposition~\ref{prop:main} is a genuinely new statement (not present in paper2a or in P4v1),
but it is \textbf{weaker than what P4 needs}: superlinear-in-every-constant is not superpolynomial,
and \S\ref{sec:compare} shows superpolynomial is the actual threshold to beat every fixed-order,
fixed-degree algebraic ODE via a Bank--Kaufman/Steinmetz-shaped bound (itself still uncited in
citable disc form, ingredient (b), untouched here).
\item Lemma~\ref{lem:many}'s use of primes $N$ with $a=\lfloor N/2\rfloor$ is one convenient
witnessing family, chosen for a clean proof that $\gcd(a,N)=1$; nothing hinges on primality beyond
that convenience, and any infinite family of coprime pairs landing in a fixed compact sub-arc away
from $\zeta=1,-1$ would do equally well.
\item The claim in \S\ref{sec:diagonal}'s proof that distinct-$\zeta_k$ pole clusters do not overlap
for $r$ close to $1$ is argued qualitatively (tie-ray convergence in both modulus and argument); it
was not turned into an explicit separation bound. For a fixed finite $K$ this is not a real risk
(compactness of the finite set $\{\zeta_1,\ldots,\zeta_k\}$ suffices), but it is the same missing
uniformity that would need to be resolved to let $K\to\infty$ with $r$, per \S\ref{sec:diagnosis}.
\item No computation (Rust, Arb, or otherwise) was run for this note beyond the arithmetic sanity
check in \S\ref{sec:local}, which reuses numbers already printed in paper2a; nothing here required
new numerics, since the content is an asymptotic inversion and a diagonal summation argument, not a
new numerical claim. If \S\ref{sec:diagnosis}'s item (3) is pursued next, that would be the first
place a fresh computation (checking whether \texttt{P1g\_lemma.tex}'s $N\to\infty$ hand argument
gives an explicit uniform bound) becomes relevant, and should be scoped as a targeted re-reading of
\texttt{P1g\_lemma.tex}/\texttt{P1h\_lemma.tex}, not a new Rust enumeration.
\end{itemize}

\end{document}
